\clearpage
\seccionInforme{6. Conclusiones}
\phantomsection\addcontentsline{toc}{section}{6. Conclusiones}

\begin{itemize}[label=\textcolor{unsared}{\textbullet}, leftmargin=*]
    \item Se desarrolló MatriLab EPIS con React Native 0.88 y Expo SDK 58: doce interfaces para tres roles que se ejecutan como aplicación nativa en Android y como aplicación web con el mismo código, usando solo componentes del núcleo del framework para organizar la interfaz.
    \item El árbol de problemas y el árbol de objetivos justificaron cada pantalla: los cinco medios fundamentales se tradujeron en trece funcionalidades y estas en las doce interfaces (Tabla \ref{tab:trazabilidad}).
    \item La investigación de componentes mostró que el layout en React Native descansa en Flexbox con valores por defecto distintos a los de la web, en listas virtualizadas para colecciones grandes y en posición absoluta y dimensiones de ventana para superposiciones y adaptación. La rejilla de laboratorios pasó de una a dos y tres columnas según el ancho sin código adicional.
    \item Las capturas de los distintos estados (errores de validación, cruces, grupos agotados, diálogos y hojas nativas, tema oscuro y orientación horizontal) confirman que cada regla de negocio tiene una respuesta visible en la interfaz.
    \item Frente a FlutterFlow (LAB 02), React Native exige escribir más, pero da control total sobre el layout y el estado y produce interfaces con los controles nativos de cada sistema; el costo principal es mantener la compatibilidad de versiones del ecosistema.
\end{itemize}

